\section{The \texorpdfstring{\Kd}{Kd} retrieval and uncertainty analysis}
\label{sec:retrieval}

Now we connect the simulator to the real Apollo data and extract \Kd. The driver
is \file{scripts/pipeline/retrieve\_kd.py}.

\subsection{The data: turning raw HFE records into ``equilibrium temperatures''}
\label{sec:data}

The bundled Apollo data lives in \file{data/apollo/depth/} as \file{.tab} text
files (one per mission per probe, e.g.\ \file{a15p1\_depth.tab}). Each row is a
timestamp, a temperature, a sensor name, and the sensor's depth in cm. The loader
is \code{load\_apollo\_hfe\_depth} (\file{lunar/validation.py}).

But the raw record drifts for years after the probes were emplaced (the drilling
disturbed the soil, which slowly re-equilibrated). We only want the final
\textbf{settled} temperature of each sensor. \code{find\_stable\_window}
(\file{lunar/apollo\_helpers.py}) scans the late part of each sensor's record and
picks the earliest window where the temperature has gone \textbf{flat} (trend
slope below \SI{0.08}{K/year}), then averages it. That average is the sensor's
\textbf{equilibrium temperature $T_{\mathrm{eq}}$}. \code{extract\_sensor\_stability}
does this for every sensor and returns depths $+$ $T_{\mathrm{eq}}$ values.

Crucially, only sensors \textbf{deeper than \code{MIN\_DEPTH\_CM = 80 cm}} are
used for the fit (the \code{deep\_mask}). Shallower sensors are contaminated by
the borestem (the hardware sticking up) and by the daily skin, so they are
excluded \emph{a priori}.

\subsection{The retrieval: sweep \texorpdfstring{\Kd}{Kd}, minimize RMSE}
\label{sec:sweep}

The core idea is a \textbf{sweep} (\code{run\_kd\_sweep\_extended}):
\begin{enumerate}[leftmargin=1.5em,itemsep=3pt]
  \item Pick a grid of candidate \Kd{} values (e.g.\ 28 values from 1 to
    \SI{15}{mW} for A15; see \code{KD\_GRIDS} in \file{config.py}).
  \item For each candidate \Kd{}, run the equilibrium solver (\code{run\_with},
    which calls \code{solve\_periodic\_equilibrium}) to get a predicted
    temperature-vs-depth profile.
  \item Read the model temperature at each real sensor depth and subtract the
    measured $T_{\mathrm{eq}}$. That difference is the \textbf{residual}.
  \item Combine the residuals into one number, the \textbf{RMSE}.
\end{enumerate}

\begin{plainenglish}
\textbf{RMSE $=$ Root Mean Square Error.} It is the standard ``how far off, on
average'' score: square every residual (so positives and negatives don't cancel
and big misses count extra), average them, take the square root. Smaller $=$
better fit. Units are kelvin here:
\[
  \mathrm{RMSE}(K_d) = \sqrt{\frac{1}{N}\sum_{i=1}^{N}
    \bigl[\,T_{\mathrm{model}}(z_i; K_d) - T_{\mathrm{eq},i}\,\bigr]^2}.
\]
\end{plainenglish}

The best-fit \Kdstar{} is the value that \textbf{minimizes} RMSE. The code even
fits a little parabola through the three points around the minimum to find the
bottom precisely (\code{kd\_star\_from\_residuals}). That's the whole retrieval:
\emph{slide the one knob \Kd{} until the model temperatures best match the
measured ones.}

\begin{figure}[h]
\centering
\includegraphics[width=0.80\textwidth]{fig_kd_retrieval.png}
\caption[The \Kd\ retrieval as an RMSE sweep]{The retrieval curve, computed from
the committed results. For each candidate deep conductivity \Kd{} we run the
simulator and score how far its temperatures land from the Apollo measurements
(RMSE). \textbf{What to notice:} each site's curve is a clean valley; its lowest
point is the retrieved \Kdstar{} (filled marker) --- \num{4.58}~\mwmk{} for
Apollo~15 and \num{8.12}~\mwmk{} for Apollo~17. \textbf{Why it matters:} the two
valleys sit at clearly different places, which is the visual heart of the
``the sites differ'' result.}
\label{fig:retrieval}
\end{figure}

Result: A15 lands at \Kdstar{}~${\approx}\num{4.58}$, A17 at
${\approx}\num{8.12}$~\mwmk{}, and switching to per-site values roughly
\textbf{halves} the Apollo-17 mismatch versus the old global \Kd{}~$=\num{3.4}$.

\subsection{Uncertainty: the bootstrap}
\label{sec:bootstrap}

A single best-fit number isn't science without an error bar. How sure are we? The
biggest uncertainty is that the \textbf{sensor depths themselves} are only known
to about $\pm\SI{2.5}{cm}$ (Nagihara~2018). The repo propagates this with a
\textbf{bootstrap} (\code{bootstrap\_kd\_with\_depth\_uncertainty}):
\begin{itemize}[leftmargin=1.4em,itemsep=3pt]
  \item Repeat 1500 times: randomly \textbf{resample} the sensors (draw with
    replacement) \emph{and} \textbf{jitter} each depth by a random
    $\pm\SI{2.5}{cm}$ wiggle, then re-find \Kdstar{}.
  \item You end up with 1500 slightly different \Kdstar{} values --- a whole
    distribution.
  \item The middle $95\%$ of that distribution is the \textbf{95\% confidence
    interval} (the $+1.33/-0.28$ style error bars).
\end{itemize}

\begin{analogy}
A bootstrap is just ``shake the inputs within their uncertainty, many times, and
watch how much the answer wobbles.'' The wobble \emph{is} the error bar. (The
bootstrap reuses cached solver outputs, so it is pure fast statistics --- no new
physics runs. It is ${\sim}65\%$ of the pipeline's runtime purely because of the
1500 repeats.)
\end{analogy}

\begin{figure}[h]
\centering
\includegraphics[width=0.92\textwidth]{fig_bootstrap.png}
\caption[Bootstrap distributions of the retrieved \Kdstar]{The bootstrap
distributions of \Kdstar{} for both sites. Each histogram is the spread of
answers you get when the sensor depths are jittered within their
$\pm\SI{2.5}{cm}$ uncertainty and the sensors are resampled, 1500 times. Dashed
lines mark the medians; shaded bands mark the 95\% confidence intervals.
\textbf{What to notice:} the two clouds barely overlap. \textbf{Why it matters:}
that separation is the quantitative statement that Apollo~17 conducts heat better
than Apollo~15 --- the contrast's confidence interval excludes zero
($p\approx0.01$).}
\label{fig:bootstrap}
\end{figure}

\subsection{Inter-site contrast and the \texorpdfstring{\Kd--\Qb}{Kd-Qb} degeneracy}
\label{sec:degeneracy}

The headline scientific claim is the \textbf{contrast}: A17~$-$~A15
${\approx}\num{3.3}$~\mwmk{}. Because we have a bootstrap distribution for
\emph{each} site, we can subtract them sample-by-sample to get a distribution for
the \textbf{difference}, and check whether it excludes zero. It does, at the
published heat fluxes: 95\% CI ${\approx}[0.4, 4.6]$, $p\approx0.01$
(\code{contrast\_bootstrap} in the results JSON).

The big caveat the paper is honest about: deep temperature depends on the
\textbf{ratio} $Q_b/K_d$, so \Kd{} and the basal flux \Qb{} are partly
\textbf{degenerate} --- if you assumed a different \Qb{}, you would retrieve a
different \Kd{}. That is why every result is stated \emph{``conditional on the
published basal heat fluxes.''} Propagating the full uncertainty in \Qb{} dilutes
the significance (the ordering probability drops to ${\sim}83\%$).

\begin{keyidea}
The honest summary: the sites differ in the combination $T(z)$ they produce;
attributing that purely to conductivity rather than basal flux rests on the
published \Qb{} values.
\end{keyidea}

Several auxiliary scripts probe robustness (held-out tests, model selection,
borestem cut, an independent Bayesian/MCMC cross-check) --- see
\autoref{sec:codetour}.

\subsection{The optional bedrock toggle (for future ice work)}
\label{sec:bedrock}

Recently added: an \textbf{optional bedrock layer} (\code{with\_bedrock} in
\file{lunar/properties.py}, configured by \code{BEDROCK} in
\file{lunar/config.py}). The standard model treats the Moon as fluffy regolith
all the way down. In reality, below ${\sim}\SI{10}{m}$ there is solid rock, which
conducts far better. \code{with\_bedrock} wraps any conductivity model so that
below \code{z\_bedrock} (\SI{10}{m}) the conductivity ramps smoothly (via a
\code{tanh}) up to a rock value (\code{K\_rock = 2.0}).

It is \textbf{OFF by default} and \textbf{does not change the published
retrieval} --- every Apollo sensor is at $\le\SI{2.4}{m}$, far above the
\SI{10}{m} transition, where enabling bedrock shifts temperatures by only
${\sim}\SI{0.05}{\kelvin}$. The tests \file{tests/test\_bedrock.py} lock in both
guarantees (off by default; sensor-depth temperatures unchanged). It exists for
\textbf{future deep-profile / ice-stability studies}, where the temperature far
below the regolith \emph{does} matter.
